%% AMS book starter — amsbook document class.
%%
%% The amsbook class itself is NOT included in this download: it ships with
%% every modern TeX distribution (and with latextemplates.com's compiler), so
%% bundling a copy would only risk pinning you to an older version than your
%% own system already has.
%%
%% Class documentation: texdoc amsbook   |   https://ctan.org/pkg/amscls
%% Licence: amsbook (part of amscls) is LPPL-1.3c, (c) American Mathematical
%% Society.
%%
%% Unlike a bare article, a book needs \frontmatter / \mainmatter /
%% \backmatter to get roman-numeral front matter, a fresh arabic page 1 at
%% the first chapter, and correctly numbered back matter — amsbook supplies
%% all three, which is the actual reason to reach for it over amsart.
\documentclass{amsbook}

\usepackage{amsmath,amssymb,amsthm}

\newtheorem{theorem}{Theorem}[chapter]
\newtheorem{lemma}[theorem]{Lemma}
\newtheorem{corollary}[theorem]{Corollary}
\newtheorem{proposition}[theorem]{Proposition}

\theoremstyle{definition}
\newtheorem{definition}[theorem]{Definition}

\theoremstyle{remark}
\newtheorem{remark}[theorem]{Remark}

\title{Foundations of Analytic Structures}
\author{Ada Lovelace}
\address{Institute of Analytical Engines, London, United Kingdom}
\email{ada@example.edu}

\date{\today}

\begin{document}

\frontmatter

\maketitle

\chapter*{Preface}
This book grew out of a lecture series on analytic structures given at the
Institute of Analytical Engines. It is intended as a graduate-level
introduction, assuming only a first course in real analysis.

\tableofcontents

\mainmatter

\chapter{Preliminaries}
\label{ch:preliminaries}

We fix notation and recall the basic definitions used throughout.

\begin{definition}[Analytic structure]
\label{def:analytic-structure}
An \emph{analytic structure} on a set $X$ is a collection $\mathcal{A}$ of
functions $X \to \mathbb{C}$ closed under addition, multiplication, and
locally uniform limits.
\end{definition}

\begin{theorem}
\label{thm:first}
Every analytic structure $\mathcal{A}$ on a compact set $X$ contains the
constant functions.
\end{theorem}

\begin{proof}
Let $c \in \mathbb{C}$ be constant. The constant function $c$ is the locally
uniform limit of the sequence $c, c, c, \dots \in \mathcal{A}$, hence lies in
$\mathcal{A}$ by closure under limits.
\end{proof}

\begin{remark}
Theorem~\ref{thm:first} is the reason we always assume $X \neq \emptyset$.
\end{remark}

See~\cite{rudin1976principles} for the classical real-analysis background,
and~\cite{ahlfors1979complex} for the complex-analytic prerequisites.

\chapter{Convergence and Compactness}
\label{ch:convergence}

\begin{proposition}
\label{prop:compact-embed}
If $X$ is compact and $\mathcal{A}$ separates points, then the evaluation map
$X \to \mathbb{C}^{\mathcal{A}}$ is an embedding.
\end{proposition}

\begin{proof}
Continuity and injectivity follow from separation of points
(Definition~\ref{def:analytic-structure}); compactness of $X$ upgrades the
continuous injection to a homeomorphism onto its image, by the standard
closed-map lemma.
\end{proof}

\begin{corollary}
Under the hypotheses of Proposition~\ref{prop:compact-embed}, $X$ is
metrizable.
\end{corollary}

\backmatter

\bibliographystyle{amsplain}
\bibliography{references}

\end{document}
